\clearpage\SourceLeaf{302}
\begin{center}
\textbf{[ZU MÖBIUS’ BARYZENTRISCHEM KALKÜL.]}

\medskip

[I.]
\end{center}

\begin{center}\textsc{Gauss} an \textsc{Schumacher}. Göttingen, 12. Mai 1843.\end{center}

Ich schulde Ihnen, mein teuerster Freund, noch meinen Dank für mehrere Ihrer Briefe und sonstige
interessante Mitteilungen. Sie dürfen aber nicht zu genau mit
mir rechnen. Seit 4---6 Wochen bin ich, anfangs durch zufällige Umstände, in einige
mathematische Spekulationen hineingezogen worden, in denen mich neue Aussichten immer wieder in
andere Richtungen gelenkt haben und ich vieles erreicht, vieles verfehlt habe. Bei einem
solchen Umhertreiben ist man (oder bin ich wenigstens) fast unfähig, sich auf irgendwelche
anderen Gegenstände einzulassen. So ist besonders in meiner Korrespondenz sehr viel liegen geblieben.
Bei jenen Spekulationen ging es größtenteils weniger um neue Dinge als um die Durchführung
meiner eigenen Methoden, zuletzt unter anderem bei mehreren Fragen zu den
Kegelschnitten. Dabei ist mir wieder eingefallen, dass ich vor einem halben
Jahrhundert, als ich zum ersten Mal \textsc{Newtons} Principia las, manches unbefriedigend
fand, namentlich auch seine an sich herrlichen Sätze über die Kegelschnitte.
Doch las ich immer mit dem Gefühl, durch das Gelernte nicht Herr der Sache
zu werden. Besonders quälte mich die Gerade, mit deren Hilfe ein Kegelschnitt
beschrieben werden kann. \textsc{Newton} löst die Aufgabe, durch 5 Punkte $A$, $B$,
$C$, $D$, $E$ einen Kegelschnitt zu legen, indem er zuerst mithilfe der vier
